\documentclass[12pt]{article}
\usepackage[utf8]{inputenc}
\usepackage[margin=1in]{geometry}
\usepackage{indentfirst}
\usepackage[osf]{mathpazo}

\date{} %date not required
\title{Information provision and externalities: The impact of genomic testing on the dairy industry} 

\author{
    Victor Funes-Leal \\ % Author 1
    Department of Agricultural Economics and Agribusiness,\\
    University of Arkansas \\
    \texttt{vf006@uark.edu}
    \and % Separator for multiple authors
    Jared Hutchins\\ % Author 2
    Department of Agricultural and Consumer Economics,\\
    University of Illinois Urbana Champaign \\
    \texttt{jhtchns2@illinois.edu}
}


\begin{document}
\maketitle

Dear Christian Langpap, 

We wish to submit an original research article entitled ``Information provision and externalities: The impact of genomic testing on the dairy industry'' for consideration by the \textbf{American Journal of Agricultural Economics}. 

This is the first research article to examine the long-term economic implications of Genomic Selection. This technique, introduced in 2009, enables dairy cattle breeders to select the best bulls (in terms of heritable traits, such as milk, protein, and fat production) to market as soon as the animal is born, rather than waiting a few years for the bulls to have a large enough number of offspring. To our knowledge, no other papers have examined this problem from an economic perspective before; similarly, no other article has attempted to estimate a causal model to properly identify the causal effects of this technology on inbreeding rates and animal-level output.

We show that genomic selection led to significant improvements in per animal productivity, at the expense of an increase in inbreeding (dairy cows are becoming more related to each other at the national level). 

Increased inbreeding implies significant present and future costs for the American and international dairy industries, due to an increased frequency of recessive traits and lower cow fertility rates, which impose an increased cost on the entire industry. We argue that this is a negative externality that will generate billions of dollars in costs in the next few years.

We, the authors, affirm that all material in this manuscript has not been published or submitted for publication elsewhere, and will not be submitted for publication elsewhere, unless rejected by the journal editor or withdrawn by the authors.

We declare that we have no material or financial interest in the subject of this research. All data is publicly available on the National Association of Animal Breeders' website. The NAAB did not provide financial or any other kind of support to the authors.

The data and code necessary to replicate the results can be shared if the article is accepted for publication.

\textbf{Word count}: 8227 words

\textbf{Table count}: 4 tables

\textbf{Figure count}: 14 figures

\end{document}